% ---
% Reemplaza líneas 1884-1906 de Docs/Technical_Report.tex
% Decisiones: D-033
% ---

\subsubsection{Visualización \textit{t-SNE}}

El algoritmo \textit{t-Distributed Stochastic Neighbor
Embedding} (t-SNE) \cite{vandermaaten2008tsne} proyecta
vectores de alta dimensionalidad a un espacio de dos
dimensiones preservando la estructura de vecindad local.
Se utiliza como herramienta de visualización cualitativa
de la separabilidad de clases en el espacio de embeddings:
si la proyección t-SNE muestra clusters compactos y bien
separados por clase, el espacio de representación tiene
estructura discriminativa.

Como complemento visual de las métricas de separabilidad del espacio de embeddings
(DB, Fisher multivariante y KTA lineal), t-SNE se emplea como herramienta de exploración cualitativa.
Para garantizar reproducibilidad, la proyección bidimensional se genera fijando
perplejidad 30, inicialización determinista basada en las dos primeras componentes principales
(PCA) y semilla aleatoria 42 (\texttt{sklearn.manifold.TSNE}). Se utiliza con fines
de soporte visual y no como criterio numérico de decisión.
